\documentclass[11pt, a4paper]{article}

\usepackage[T1]{fontenc}
\usepackage[utf8]{inputenc}
\usepackage{tgpagella}
\usepackage[spanish, es-noshorthands]{babel}
\usepackage{amsmath, amssymb, amsfonts}
\usepackage{booktabs}
\usepackage{tabularx}
\usepackage[table, xcdraw]{xcolor}
\usepackage[most]{tcolorbox}
\usepackage[margin=2.5cm]{geometry}
\usepackage{fancyhdr}

\pagestyle{fancy}
\fancyhf{}
% Aumentamos el headheight por si acaso y dejamos el centro vacío para evitar solapamientos
\setlength{\headheight}{16pt}
\lhead{\textbf{UCB Sede Tarija} | Ing. Mecatrónica}
\chead{}
\rhead{IMT-121 Circuitos Electrónicos I | Defensa Hito 1}
\lfoot{Prof: M.Sc. Ing. Bernardo Quiroga Turdera}
\rfoot{Página \thepage}
\renewcommand{\headrulewidth}{0.4pt}

\definecolor{ucbblue}{RGB}{0, 51, 102}

\begin{document}

\begin{tcolorbox}[colback=ucbblue!5!white, colframe=ucbblue, title=\textbf{Hoja de Ruta y Diagnóstico: Defensa del Hito 1}]
    \centering \large \textbf{Validación de Modelado Matricial y Simulación Computacional}
\end{tcolorbox}

\section*{Secuencia Recomendada por Equipo (55-65 min total aprox.)}
\begin{enumerate}
    \item \textbf{30--60 s:} Propósito del circuito y variables.
    \item \textbf{1--2 min:} Seleccionar una ecuación y reconstruirla físicamente.
    \item \textbf{1--2 min:} Mapear ecuación $\leftrightarrow$ fila/coeficientes de $A$.
    \item \textbf{1--2 min (por estudiante):} Pregunta individual de una categoría distinta (muestreo estratégico).
    \item \textbf{1--2 min:} Verificación en Python y orden de variables en el vector $x$.
    \item \textbf{1--2 min:} Correspondencia de magnitudes en LTSpice.
    \item \textbf{1 min:} Discrepancia/validación y función del Hito 1 en el Proyecto Final de Ingeniería (PFI).
\end{enumerate}

\section*{Banco Compacto de Preguntas Equivalentes}
\begin{tcolorbox}[colback=white, colframe=black, boxrule=0.5pt, arc=2mm]
\textbf{Modelo / Ecuaciones:}
\begin{itemize}\setlength\itemsep{-0.2em}
    \item Explique una ecuación directamente desde el circuito.
    \item ¿De dónde sale físicamente este coeficiente en la ecuación?
    \item ¿Por qué este término tiene un signo negativo/positivo según su convención?
\end{itemize}
\textbf{Matriz ($A\cdot x=b$):}
\begin{itemize}\setlength\itemsep{-0.2em}
    \item ¿Qué representa físicamente esta fila de la matriz $A$?
    \item ¿Cuál es el orden exacto de las variables en su vector $x$?
    \item Si se cambia el valor de esta fuente (o resistencia), ¿qué cambia en el sistema matricial?
\end{itemize}
\textbf{Python / NumPy:}
\begin{itemize}\setlength\itemsep{-0.2em}
    \item ¿Qué está resolviendo matemáticamente la función \texttt{solve}?
    \item ¿Dónde aparece esta variable específica en su vector de resultados $x$?
    \item ¿Qué verifica exactamente la operación $A @ x \approx b$?
\end{itemize}
\textbf{LTSpice:}
\begin{itemize}\setlength\itemsep{-0.2em}
    \item ¿Qué salida (.op) de LTSpice corresponde a esta variable analítica?
    \item ¿La corriente medida en simulación tiene el mismo sentido de referencia que su modelo?
    \item ¿Esta corriente de malla se obtiene directamente o requiere derivarse (restarse) de otras corrientes?
\end{itemize}
\textbf{Validación y PFI:}
\begin{itemize}\setlength\itemsep{-0.2em}
    \item ¿Qué dos magnitudes (y unidades) está comparando exactamente en su cálculo de error?
    \item ¿Cómo explica esta discrepancia (qué hipótesis tiene)?
    \item ¿Qué demuestra la coincidencia numérica y qué \textit{no} demuestra?
    \item ¿Qué aporta el modelo validado en este Hito 1 a los siguientes pasos del PFI?
\end{itemize}
\end{tcolorbox}
\newpage

\section*{Registro de Diagnóstico Individual (Rápido)}
\vspace{0.2cm}

\noindent
\renewcommand{\arraystretch}{1.5}
\begin{tabularx}{\textwidth}{|>{\columncolor{ucbblue!5}\bfseries}l|X|}
\hline
Estudiante 1: & \\
\hline
Fortalezas: & $\square$ Modelo \quad $\square$ Ecuaciones \quad $\square$ Matriz \quad $\square$ Python \quad $\square$ LTSpice \quad $\square$ Interp. \\
\hline
Brechas: & $\square$ KCL-KVL \quad $\square$ Signos \quad $\square$ Álgebra \quad $\square$ Python \quad $\square$ LTSpice \quad $\square$ Val./PFI \\
\hline
Evidencia breve: & \\
\hline
Autonomía: & $\square$ Sólida \quad $\square$ Parcial \quad $\square$ Insuficiente \\
\hline
\end{tabularx}

\vspace{0.4cm}

\noindent
\begin{tabularx}{\textwidth}{|>{\columncolor{ucbblue!5}\bfseries}l|X|}
\hline
Estudiante 2: & \\
\hline
Fortalezas: & $\square$ Modelo \quad $\square$ Ecuaciones \quad $\square$ Matriz \quad $\square$ Python \quad $\square$ LTSpice \quad $\square$ Interp. \\
\hline
Brechas: & $\square$ KCL-KVL \quad $\square$ Signos \quad $\square$ Álgebra \quad $\square$ Python \quad $\square$ LTSpice \quad $\square$ Val./PFI \\
\hline
Evidencia breve: & \\
\hline
Autonomía: & $\square$ Sólida \quad $\square$ Parcial \quad $\square$ Insuficiente \\
\hline
\end{tabularx}

\vspace{0.4cm}

\noindent
\begin{tabularx}{\textwidth}{|>{\columncolor{ucbblue!5}\bfseries}l|X|}
\hline
Estudiante 3: & \\
\hline
Fortalezas: & $\square$ Modelo \quad $\square$ Ecuaciones \quad $\square$ Matriz \quad $\square$ Python \quad $\square$ LTSpice \quad $\square$ Interp. \\
\hline
Brechas: & $\square$ KCL-KVL \quad $\square$ Signos \quad $\square$ Álgebra \quad $\square$ Python \quad $\square$ LTSpice \quad $\square$ Val./PFI \\
\hline
Evidencia breve: & \\
\hline
Autonomía: & $\square$ Sólida \quad $\square$ Parcial \quad $\square$ Insuficiente \\
\hline
\end{tabularx}

\end{document}
